\section{从平均变化率到导数}
\label{sec:02-Calculus/03-Differentiation/01}

\subsection{速度表上的数字是怎么来的}

一辆车在高速公路上行驶。你盯着速度表，指针稳在 100 km/h。但速度表并不是在测量一段路程再除以时间——它给出的是 ``此刻'' 的速度。

可 ``此刻'' 没有长度。没有位移，没有时间间隔，速度该怎么定义？你真正能量的是从 \(t\) 到 \(t+h\) 这段时间内的平均速度：

\[
\bar v=\frac{s(t+h)-s(t)}{h}.
\]

但 \(h=1\) 秒的平均速度和 \(h=0.01\) 秒的平均速度可能完全不同。哪一个才算 ``时刻 \(t\) 的速度''？

假设位置函数为 \(s(t)=t^2\)（单位：米，\(t\) 以秒计）。在 \(t=3\) 这个时刻，取不同长度的间隔：

\[
\begin{aligned}
h=0.5&\colon \frac{(3.5)^2-3^2}{0.5}=\frac{12.25-9}{0.5}=6.5,\\
h=0.1&\colon \frac{(3.1)^2-3^2}{0.1}=\frac{9.61-9}{0.1}=6.1,\\
h=0.01&\colon \frac{(3.01)^2-3^2}{0.01}=\frac{9.0601-9}{0.01}=6.01,\\
h=0.001&\colon \frac{(3.001)^2-3^2}{0.001}=6.001.
\end{aligned}
\]

数值在向 6 靠近。\(h\) 每缩小一个数量级，平均速度就向 6 多逼近一位小数。这个趋势是清楚的。

但你不能直接令 \(h=0\)——那会得到 \(\frac{0}{0}\)，毫无意义。这里的困境与极限计算时一模一样：你想要某个量在 \(h\to0\) 时的行为，却不能直接在表达式里写 \(h=0\)。

\subsection{几何那边也在发生同类事件}

不只是速度。画一条曲线 \(y=f(x)\)，你想知道在点 \((a,f(a))\) 处曲线 ``朝哪个方向走''。

你能画的是割线——连接 \((a,f(a))\) 和附近某点 \((a+h,f(a+h))\) 的直线。割线的斜率是

\[
\frac{f(a+h)-f(a)}{h}.
\]

取 \(f(x)=x^2\)，\(a=1\)。与刚才完全相同的代数结构：

\[
\begin{aligned}
h=0.5&\colon \frac{(1.5)^2-1^2}{0.5}=2.5,\\
h=0.1&\colon \frac{1.21-1}{0.1}=2.1,\\
h=0.01&\colon 2.01.
\end{aligned}
\]

当 \(h\) 越来越小，割线围绕 \((1,1)\) 旋转，斜率越来越靠近 2。极限位置上那条 ``刚好擦过曲线'' 的直线，就是切线。

物理的瞬时速度和几何的切线斜率，在数学上是同一个结构：差商的极限。一旦你解决了这个极限问题，两边同时受益。

\subsection{极限接管了除法的烂摊子}

出路你已经见过了：用极限。先对差商做代数化简（在 \(h\neq0\) 的穿孔邻域内合法），再取极限：

\[
\frac{(t+h)^2-t^2}{h}
=\frac{t^2+2th+h^2-t^2}{h}
=\frac{2th+h^2}{h}
=2t+h,\qquad h\neq0.
\]

当 \(h\to0\)，\(2t+h\to2t\)。所以时刻 \(t\) 的瞬时速度为 \(2t\)。在 \(t=3\) 处就是 6 m/s，与数值试探吻合。

这里有一个微妙但重要的操作顺序：先在 \(h\neq0\) 的条件下约去 \(h\)，再让 \(h\to0\) 取极限。整个过程没有哪一步把 \(h\) 直接设成了零。约分用的是代数恒等式，极限用的是前面建立的 \(\varepsilon\)--\(\delta\) 机制或者运算法则。两件事分开做，互不侵权。

单位也自然携带过来：若位置以米、时间以秒计，差商的单位是米/秒，导数的单位同样是米/秒。速度表上的数字就是这样从位置变化中提取出来的——不是魔法，是极限。

\subsection{导数：切线、速度，同一个定义}

函数 \(f\) 在 \(a\) 的导数定义为

\[
f'(a)=\lim_{h\to0}\frac{f(a+h)-f(a)}{h},
\]

若这个有限极限存在。等价地，也可写成

\[
f'(a)=\lim_{x\to a}\frac{f(x)-f(a)}{x-a}.
\]

两个形式表达同一件事：函数在 \(a\) 附近的平均变化率，当观测窗口无限缩窄时，是否稳定到一个确定的值。

记号上，\(f'(a)\) 是拉格朗日记号，\(\frac{dy}{dx}\) 是莱布尼茨记号。两者等价，但莱布尼茨记号在链式法则和多变量情形中更直观——它保留了 ``变化率相比什么在变化'' 的信息。

切线方程写为

\[
y=f(a)+f'(a)(x-a).
\]

这不只是 ``在 \(a\) 处接触图像''——它是函数在 \(a\) 附近的一阶最佳线性近似。在 \(a\) 旁边取一点 \(x\)，用切线值 \(f(a)+f'(a)(x-a)\) 代替真实函数值 \(f(x)\)，误差随 \(|x-a|\) 的缩小而比线性更快地趋于零。换句话：当你看得足够近，曲线和切线几乎无法分辨。

以 \(f(x)=x^2\) 在 \(a=1\) 为例，\(f'(1)=2\)，切线方程为 \(y=1+2(x-1)=2x-1\)。在 \(x=1.01\) 处，真实值 \(1.0201\)，切线近似值 \(1.02\)，误差仅 \(0.0001\)。

\subsection{可导一定连续，但反过来不一定}

若 \(f'(a)\) 存在，那么

\[
f(a+h)-f(a)
=h\cdot\frac{f(a+h)-f(a)}{h}.
\]

当 \(h\to0\)，前一因子趋于 0，后一因子趋于有限的 \(f'(a)\)，乘积趋于 0。这意味着 \(f(a+h)\to f(a)\)——函数在 \(a\) 处连续。可导性自动携带了连续性：切线存在的前提是图像在这一点没有断开。

反过来呢？连续不能保证可导。

最经典的例子是 \(f(x)=|x|\) 在 \(x=0\)。函数在 0 处连续：\(|0|=0\)，左右都向 0 靠近。但看差商：

\[
\lim_{h\to0^+}\frac{|h|-0}{h}=\frac{h}{h}=1,\qquad
\lim_{h\to0^-}\frac{|h|-0}{h}=\frac{-h}{h}=-1.
\]

从右边逼近，割线斜率是 1；从左边逼近，割线斜率是 \(-1\)。两条割线旋转到不同的极限位置——函数在此处有一个尖角，不存在唯一的切线方向。左右导数不相等，导数不存在。

再看 \(f(x)=x^{1/3}\) 在 \(x=0\)。它处处连续，但在 0 处的差商是 \(h^{1/3}/h=h^{-2/3}\)，当 \(h\to0\) 时无界——切线是竖直的，斜率不存在为有限数。

连续只要求点值接上，导数还要求局部方向稳定成一条有限斜率的直线。一个只关心 ``有没有断''，另一个关心 ``拐得顺不顺''。这个区别在实际中很重要：优化算法依赖导数指引方向，\(|x|\) 这种尖角处，算法不知道该往左还是往右走。

\subsection{导数符号告诉你什么——以及不告诉你什么}

\(f'(a)>0\) 表示函数在 \(a\) 的局部线性趋势向上：在足够小的邻域内，\(x>a\) 的点图像高于 \(f(a)\)，\(x<a\) 的点低于 \(f(a)\)。\(f'(a)<0\) 表示向下。\(f'(a)=0\) 表示切线水平。

但导数为零不意味着极大值或极小值。\(f(x)=x^3\) 在 \(x=0\) 处的导数为零——切线确实是水平的——可函数仍在穿过原点继续增长：\(x<0\) 时 \(x^3<0\)，\(x>0\) 时 \(x^3>0\)。水平切线只说明一阶信息在此处 ``熄火''，极值判断还需要更高阶的信息：二阶导数可以区分 ``穿过'' 还是 ``折返''。

同样，某点导数为正是局部信息。要断言整个区间上的单调递增，还需要导数在整个区间内保持非负，以及中值定理把逐点信息串联成区间结论。一个点的斜率管不了远处的事。比如 \(f(x)=x+\sin(10x)\)，在 \(x=0\) 处导数是正数，但函数的图像充满剧烈振荡，在任意更长的区间上并非单调。

\subsection{从差商到导数：你得到了什么工具}

总结一下这条推导链给了你什么。给定一个函数 \(f\)，求它在某一点的导数，标准动作是三步：

（1）写出差商 \(\frac{f(a+h)-f(a)}{h}\)；（2）在 \(h\neq0\) 的条件下化简；（3）令 \(h\to0\) 取极限。

这三步的背后，是一次认知升级：你不再把变化率理解为 ``一段区间上的平均''，而是 ``一个点的瞬时趋势''。这个升级的代价是：你必须接受极限作为合法的数学操作，而不能期待每次都能直接代入 \(h=0\)。

导数给你的不是一个近似的速度，而是精确的切线斜率。但精确不等于无限制——它只在函数足够光滑的点上有定义，而且它只告诉你一阶信息。弯曲、拐点、全局走势，还需要后续的求导规则、二阶导数和中值定理来逐步揭开。
